% !TITLE 月夜
% !AUTHOR 杜甫
% !DYNASTY 唐
% !ORDER 23

\begin{poem}
今夜鄜州月，闺中只独看。
遥怜小儿女，未解忆长安。
香雾云鬟湿，清辉玉臂寒。
何时倚虚幌，双照泪痕干。
\end{poem}

\begin{appreciation}
这是杜甫被困长安时写给妻子的诗。至德元年（756年），安史之乱爆发，杜甫把家人安顿在鄜州（今陕西富县），自己则去投奔新即位的唐肃宗，不料中途被叛军俘获，押回长安。

今夜鄜州月，闺中只独看——今夜鄜州的月亮，妻子只能在闺中独自观看。杜甫不写自己看月，而写妻子看月。这种从对方角度着笔的写法，让思念变得更加深切。只独看三个字，藏着多少孤寂。

遥怜小儿女，未解忆长安——可怜我那年幼的儿女，还不懂得思念长安的父亲。未解是最让人心酸的：孩子太小了，连想念都不会。杜甫想念他们，但他们还不知道想念是什么。这种不对等的思念，让父亲的心更加疼痛。

香雾云鬟湿，清辉玉臂寒——夜雾打湿了她的发髻，月光让她的手臂感到寒冷。云鬟是古代女子的发式，玉臂是对妻子手臂的美称。杜甫想象着妻子在月下站立太久，雾气沾湿了头发，月光冻冷了手臂。这些细节如此具体，仿佛他亲眼看见了一般。这就是杜甫的想象力：不是天马行空，而是设身处地。

何时倚虚幌，双照泪痕干——什么时候才能一起倚在薄帷前，让月光照着我们两个人，把泪痕照干？虚幌就是轻薄的帷幔。双照是两个人一起被月光照耀，对应前面的独看。泪痕干意味着流泪已经结束，意味着团聚已经到来。但何时两个字，说明这只是一个愿望，一个不知道能不能实现的愿望。
\end{appreciation}
